%% ma: L12-B06-0087
%% lop: 12
%% chuong: 2
%% bai: 6
%% dang: 12.6.23
%% loai: TN4
%% muc: NB
%% dap_an: "D"
%% nguon: "Ảnh giáo viên gửi 05/11/2026 (M11): Kiểm tra giữa HK1 Toán 12 (2026-2027), THPT Hồ Thị Bi (TP.HCM), mã đề 1133, Phần I Câu 14"
%% kien_thuc: "Bài 6 (vectơ cùng hướng, ngược hướng, trung điểm đoạn thẳng)"
%% trang_thai: cho-duyet
%% ly_do: "Dạng mới 12.6.23 đang đề xuất, chờ giáo viên duyệt"
%% ghi_chu: "Đề gốc không kèm đáp án; đã giải độc lập."
\begin{ex}
Cho tứ diện $ABCD$, gọi $I$ là trung điểm của cạnh $AB$ (tham khảo hình vẽ). Vectơ $\overrightarrow{IA}$ ngược hướng với vectơ nào dưới đây?
\begin{center}
\begin{tikzpicture}[scale=.9]
  \coordinate (A) at (0,2.3); \coordinate (B) at (-1.8,0); \coordinate (D) at (1.8,0); \coordinate (C) at (0,-1.1);
  \coordinate (I) at ($(A)!.5!(B)$);
  \draw (A)--(B) (A)--(D) (A)--(C) (B)--(C) (C)--(D);
  \draw[khuat] (B)--(D);
  \fill (I) circle(1pt);
  \path (A) node[above]{$A$} (B) node[left]{$B$} (D) node[right]{$D$} (C) node[below]{$C$} (I) node[left]{$I$};
\end{tikzpicture}
\end{center}
\choice{$\overrightarrow{BA}$}{$\overrightarrow{BI}$}{$\overrightarrow{CD}$}{\True $\overrightarrow{AB}$}
\loigiai{$I$ nằm giữa $A$ và $B$ nên $\overrightarrow{IA}$ cùng hướng với $\overrightarrow{BI}$ và $\overrightarrow{BA}$, ngược hướng với $\overrightarrow{AB}$.
Vectơ $\overrightarrow{CD}$ không cùng phương với $\overrightarrow{IA}$ nên không ngược hướng.}
\end{ex}
